\section{System and Failure Model}
\label{sec:model}

This section condenses a formal model held in the
artifact, which states every claim below with a citation to the enforcing
code by file and line and is machine-checked in CI: a script fails the build if a
cited path stops existing or a cited line falls outside its file. That check
proves the citation ranges are valid, not that the semantics are right. The
semantic evidence is the per-property regression tests named in
\cref{sec:protocol}.

\subsection{Principals and store}

A \emph{runner} performs at most one external mutation per invocation. A
\emph{recovery service} is read-only with respect to the external system and
writes only coordination state. A \emph{connector} is the only component that
touches the external API. The store is a single self-hosted Redis~7.2 instance
with the append-only file enabled and RDB snapshots disabled. There is no
replica, no Sentinel and no Cluster.

Redis executes a Lua script without interleaving, so a script is a single
serialization point. This is the only atomicity primitive the protocol has and
every invariant below is expressed inside one.

\subsection{Clocks}
\label{sec:model-clocks}

There is no synchrony assumption beyond lease TTLs. All ordering decisions
between principals are taken against the Redis server clock; no cross-worker
clock comparison occurs. Process-local monotonic time is used only for
heartbeat intervals and watchdogs, never for an authoritative decision. The
lease budget $T_{\text{client}} \le T_{\text{lock}} - B$ with $B \ge 15$\,s is
enforced at configuration time. It is a timing budget rather than a
guarantee: a scheduling gap remains between a successful preflight and
transmission, and we declare it as such.

\subsection{The external endpoint}
\label{sec:model-endpoint}

The endpoint is non-idempotent and non-cooperative: it need not accept an
idempotency key, and its response may be absent, late, or self-contradictory.
Two things about it are modeled explicitly, and the second is the axis the
evaluation turns on.

\paragraph{Evidence at mutation time} Ambiguity is the default here, not a
special case. The runner accepts only policy-declared
success and failure values, required non-empty and disjoint; any other value,
and any exception whatever, is coerced to ambiguity.

\paragraph{Reconciliation capability} What the endpoint can be asked after the
fact is a declared, typed property with three values
(\cref{tab:capabilities}). The declaration fails closed: a missing, null or
unrecognized capability is treated as the weakest, not as the strongest.
Classification is total. In particular, ``the endpoint reports the mutation
absent'' combined with a positive-only capability is a \emph{named contract
violation} that resolves to ambiguity, not a fall-through to ``it did not
happen''.

\begin{table}[t]
\centering
\caption{Endpoint reconciliation capability. The right-hand column is the
ceiling on what any protocol can conclude, and \cref{sec:eval-rq1} shows
AEP's declared-ambiguity rate tracking it.}
\label{tab:capabilities}
\small
\begin{tabular}{@{}p{0.30\columnwidth}p{0.60\columnwidth}@{}}
\toprule
Capability & Conclusions the endpoint can support \\
\midrule
\textsc{authoritative} \textsc{readback} &
  applied, not-applied, unknown, conflict. Absence \emph{proves} the mutation
  did not occur. \\
\textsc{positive-only} \textsc{readback} &
  applied, unknown, conflict. Presence can be confirmed; absence proves
  nothing. \\
\textsc{no-readback} &
  none. The endpoint cannot be queried at all. \\
\bottomrule
\end{tabular}
\end{table}

\subsection{Failure model}

The model is crash-and-delay, not Byzantine. Redis is trusted to execute
scripts atomically
and honor expiry; the provider is trusted only insofar as its declared
evidence is truthful.

\begin{itemize}
\item[\textbf{F1}] \textbf{Worker crash at any instruction boundary.} Six named
      crash points around the external call
      (\cref{tab:crashpoints}). In the evaluation these are OS-level
      \texttt{SIGKILL} of a separate worker process, not in-process
      exceptions.
\item[\textbf{F2}] \textbf{Partition between a worker and Redis.} Any Redis
      command may raise; every raise forbids progress rather than permitting a
      guess. A partition outlasting the lease is observationally equal to F5.
\item[\textbf{F3}] \textbf{Redis restart with AOF replay.} Up to roughly one
      second of acknowledged-but-unsynced writes may be absent under
      \texttt{appendfsync everysec}~\cite{redis-persistence}.
      \Cref{sec:eval-durability} reports two probes that sharpen this
      considerably. The set of faults that can actually realize that loss is
      narrower than the documentation's wording suggests, and a process kill
      is not among them; a host-level write loss is, and realizes the loss
      every time.
\item[\textbf{F4}] \textbf{Delayed or duplicated external responses.} A late
      effect cannot be recalled; recovery waits out a reconciliation delay
      before drawing any conclusion, and never re-dispatches.
\item[\textbf{F5}] \textbf{Worker pause past lease expiry} (GC or VM stall),
      handled by an atomic pre-dispatch preflight, cancellation of the
      protected task on ownership loss, and CAS rejection of the late write.
\end{itemize}

Out of model: Byzantine Redis; a provider whose declared evidence contradicts
its own effects; loss of the Redis host once the fsync acknowledgment has been
returned; adversarial writers with direct socket access; compromise of vault
keys.

\begin{table}[t]
\centering
\caption{The six crash points, and what exists in the world at each. The
right-hand column is why the declared-ambiguity rate in \cref{sec:eval-rq1} varies
so sharply across them.}
\label{tab:crashpoints}
\small
\renewcommand{\arraystretch}{1.08}
\begin{tabularx}{\columnwidth}{@{}>{\raggedright\arraybackslash}p{0.46\columnwidth}>{\raggedright\arraybackslash}X@{}}
\toprule
Crash point & State when the worker dies \\
\midrule
\path{before_intent_write}            & no record, no effect \\
\path{after_intent_before_barrier}     & record in memory, maybe not on disk; no effect \\
\path{after_barrier_before_dispatch}   & record durable; \textbf{no effect} \\
\path{mid_dispatch}                    & record durable; effect \textbf{likely} \\
\path{after_response_before_resolution} & record durable; effect known to the dead process only \\
\path{after_resolution_before_barrier} & outcome written, not yet acknowledged \\
\bottomrule
\end{tabularx}
\end{table}

\subsection{Non-claims}
\label{sec:nonclaims}

We state these because reviewers are right to look for them, and because
several are structural rather than incidental.

\begin{table}[t]
\centering
\caption{Non-claims, with the reason each is out of scope.}
\label{tab:nonclaims}
\small
\begin{tabular}{@{}p{0.36\columnwidth}p{0.54\columnwidth}@{}}
\toprule
Non-claim & Why \\
\midrule
Exactly-once external execution
  & Redis and a legacy provider share no transaction coordinator. The
    protocol converts this into declared ambiguity. \\
Duplicate \emph{prevention}
  & AEP prevents \emph{automatic} duplicates. An effect the provider already
    accepted cannot be recalled. \\
HA, consensus, split-brain immunity
  & Single instance; the lock is a lease~\cite{gray1989leases}. This is also
    why the AOF-rewind residual below has no local fix. \\
A hardened trust domain
  & No Redis ACL; any principal reaching the socket can write the state key
    directly. \\
Durability beyond one local AOF fsync
  & The barrier requires one \emph{local} fsync~\cite{redis-waitaof}; it is
    explicitly not a substitute for replication. \\
Operator alerting
  & Not implemented. ``Escalates'' means \emph{reaches a terminal automated
    state and pauses the execution}, nothing more. \\
A verified implementation
  & The properties are model-checked (\cref{sec:modelchecking}); the Python is
    not. Refinement from the specification to the code is neither proven nor
    claimed. \\
\bottomrule
\end{tabular}
\end{table}
